\documentclass[11pt,a4paper]{article}

\usepackage[T1]{fontenc}
\usepackage[utf8]{inputenc}
\usepackage{lmodern}
\usepackage{amsmath,amssymb,amsthm}
\usepackage{mathtools}
\usepackage{booktabs}
\usepackage{tabularx}
\usepackage{enumerate}
\usepackage[margin=2.6cm]{geometry}
\usepackage[hidelinks]{hyperref}

% ---------------------------------------------------------------- notation
\newcommand{\sep}{\mathord{/}}              % the separator symbol
\newcommand{\eps}{\varepsilon}
\newcommand{\N}{\mathsf{N}}                 % the names
\newcommand{\Nat}{\mathbb{N}}
\newcommand{\U}{U}
\newcommand{\F}{F}
\newcommand{\I}{I}
\newcommand{\C}{\mathrel{C}}
\newcommand{\nm}{\operatorname{name}}
\newcommand{\Key}{\operatorname{Key}}
\newcommand{\seqof}{\operatorname{seq}}
\newcommand{\depth}{\operatorname{depth}}
\newcommand{\join}{\operatorname{join}}
\newcommand{\cat}{\mathbin{{}^\frown}}      % concatenation of sequences
\newcommand{\sq}[1]{\langle #1\rangle}      % a finite sequence
\newcommand{\pre}[2]{#1{\restriction}#2}     % prefix of length
\newcommand{\lean}[1]{\texttt{\small #1}}
\newcommand{\leanref}[1]{{\footnotesize\textup{[Lean:~\lean{#1}]}}\hskip 0.5em plus 0.2em\relax\ignorespaces}

\theoremstyle{plain}
\newtheorem{theorem}{Theorem}[section]
\newtheorem{lemma}[theorem]{Lemma}
\newtheorem{proposition}[theorem]{Proposition}
\newtheorem{corollary}[theorem]{Corollary}
\theoremstyle{definition}
\newtheorem{definition}[theorem]{Definition}
\newtheorem{axiom}{Axiom}
\renewcommand{\theaxiom}{AX-\arabic{axiom}}
\theoremstyle{remark}
\newtheorem{remark}[theorem]{Remark}

\title{The Store Model as a Set of Axioms\\[0.3em]
\large A human-readable account of the Lean~4 proofs and the Alloy checks}
\author{remote-store, BK-435 (exposition)}
\date{Written 2026-10-14 against branch \texttt{store-model-axioms} at \texttt{21a0e534b}\\
(Lean file and spec~051 unchanged at \texttt{bffcee464})}

\begin{document}
\maketitle

\begin{abstract}
A store is a rooted tree of folders and items, each addressed by one Key.
Spec~051 states that tree as eight axioms, and states the consequences a
caller relies on: one entity per Key, the shape of a Key, and closure of Keys
under ancestors. This note proves these in ordinary mathematical prose,
following the Lean~4 formalization. It also proves the formalization's
further result: up to isomorphism, the models are the non-empty prefix-closed sets of
name sequences in which the items are marked, every item a non-empty leaf.
No result proved in Lean assumes the store is finite.
A last section explains what the bounded Alloy checks add.
\end{abstract}

\paragraph{Status of this document.}
This is an \emph{exposition}, not a source of truth. The authority for the
statements is spec~051, and the machine-checked authority for the proofs is
the Lean file. Each result proved in Lean names, after its heading, the Lean
declaration it restates. One result, Proposition~\ref{prop:finitehalf}, is
checked only in Alloy and says so. Where this text and the Lean file disagree, the Lean file is right
and this text is wrong. The proofs follow the Lean proofs step by step, with
the bookkeeping that a proof assistant needs (index arithmetic, coercions,
choice of representatives) told in words.

\tableofcontents

% =======================================================================
\section{Notation and the axioms}
\label{sec:axioms}

\paragraph{Words.}
Fix an alphabet $\Sigma$ and a distinguished symbol $\sep\in\Sigma$, the
\emph{separator}. $\Sigma^*$ is the set of finite words over $\Sigma$, and
$\eps$ is the empty word. Words are compared exactly, symbol by symbol.
Juxtaposition $uv$ denotes the concatenation of words.
The \emph{names} are
\[
  \N \;=\; \{\, w\in\Sigma^* : w\neq\eps \text{ and } \sep \text{ does not occur in } w \,\}.
\]
Spec~051 asks $\Sigma$ to contain at least one symbol besides $\sep$. The proofs
below need that only in Theorem~\ref{thm:infinite}; elsewhere $\Sigma$ may be
$\{\sep\}$, in which case $\N$ is empty.

\paragraph{Sequences.}
A finite sequence of names is written $\sq{w_1,\dots,w_n}$, the empty
sequence is $\sq{}$, $s\cat t$ is the concatenation of sequences, $|s|$ is the
length, and $\pre{s}{i}$ is the prefix of $s$ of length $i$ (for $i\le|s|$).
Sequences of names and words are different things: a sequence of names is a
list of words, and the bridge between them is the join of
Definition~\ref{def:join}.

\paragraph{Structures.}
A \emph{structure} is a tuple $(\U,\F,\I,r,C,\nm)$ where $\U$ is any set (the
universe of \emph{entities}), $\F,\I\subseteq\U$ (\emph{folders} and
\emph{items}), $r\in\U$ (the \emph{root}), $C\subseteq\U\times\U$ is a binary
relation (\emph{containment}; $p\C x$ reads ``$p$ contains $x$''), and
$\nm\colon\U\to\Sigma^*$. Nothing is assumed about the size of $\U$.
\leanref{Structure}

\begin{definition}[Chain]\label{def:chain}
\leanref{IsChain}
For $x\in\U$ and $n\in\Nat$, a \emph{chain to $x$ of length $n$} is a
sequence $x_0,x_1,\dots,x_n$ of entities with $x_0=r$, $x_n=x$, and
$x_i\C x_{i+1}$ for every $i<n$.
\end{definition}

\paragraph{The axioms.}
A \emph{model} is a structure that satisfies the following. The axiom
numbers are spec~051's; TREE-00$k$ is AX-$k$.
\leanref{Axioms}

\begin{axiom}[partition]\label{ax:1} $\U=\F\cup\I$ and $\F\cap\I=\emptyset$.\end{axiom}
\begin{axiom}[root]\label{ax:2} $r\in\F$.\end{axiom}
\begin{axiom}[items are leaves]\label{ax:3} $C\subseteq\F\times\U$: if $p\C x$ then $p\in\F$.\end{axiom}
\begin{axiom}[at most one container]\label{ax:4} If $p\C x$ and $q\C x$ then $p=q$.\end{axiom}
\begin{axiom}[root is topmost]\label{ax:5} There is no $p$ with $p\C r$.\end{axiom}
\begin{axiom}[reachability]\label{ax:6} Every $x\in\U$ has a chain (of some length $n\ge0$).\end{axiom}
\begin{axiom}[names]\label{ax:7} $\nm(r)=\eps$, and $\nm(x)\in\N$ for every $x\neq r$.\end{axiom}
\begin{axiom}[sibling names]\label{ax:8} If $p\C x$, $p\C y$ and $\nm(x)=\nm(y)$, then $x=y$.\end{axiom}

AX-1 and AX-7 are conjunctions. The Lean file has one field per clause:
AX-1 \emph{disjoint} and \emph{covering}, and AX-7 \emph{root},
\emph{non-empty} and \emph{no~$\sep$}. Section~\ref{sec:alloy} refutes each clause
separately. AX-4 asserts only uniqueness. That every non-root \emph{has} a
container is a theorem (D1). No finiteness axiom is included.

% =======================================================================
\section{Words: the join of a sequence of names}
\label{sec:words}

\begin{definition}[Join]\label{def:join}
\leanref{join}
For a sequence of words, $\join\sq{}=\eps$, $\join\sq{w}=w$, and
\[
  \join\sq{w,v,\dots} \;=\; w\sep\join\sq{v,\dots}.
\]
So $\join\sq{w_1,\dots,w_n}=w_1\sep w_2\sep\cdots\sep w_n$.
\end{definition}

Everything about Keys that is not about trees is a fact about this
function. The central fact is that $\join$ loses no information on sequences
of names (Lemma~\ref{lem:joininj}). Its engine is the following
observation about where the first separator is.

\begin{lemma}[First separator]\label{lem:splitfirst}
\leanref{split\_first}
Let $w,a,\rho,J\in\Sigma^*$ with $\sep$ not occurring in $w$, and suppose
$a\sep\rho = w\sep J$. Then either $a=w$ and $\rho=J$, or there is $a'$ with
$a=w\sep a'$ and $a'\sep\rho=J$.
\end{lemma}

\begin{proof}
Induction on the length of $w$.
If $w=\eps$, the equation reads $a\sep\rho=\sep J$. If $a=\eps$, then
$\rho=J$ and the first alternative holds. If $a=b\,a''$ for a symbol $b$, then
comparing first symbols gives $b=\sep$ and $a''\sep\rho=J$, so the second
alternative holds with $a'=a''$.

If $w=c\,w'$ for a symbol $c$, then $c\neq\sep$ and $\sep\notin w'$. If
$a=\eps$, the equation reads $\sep\rho=c\,w'\sep J$, which forces $c=\sep$,
a contradiction. So $a=b\,a''$, and comparing first symbols gives $b=c$ and
$a''\sep\rho=w'\sep J$. By induction, either $a''=w'$ and $\rho=J$, giving
$a=c\,w'=w$, or $a''=w'\sep a'$ with $a'\sep\rho=J$, giving
$a=c\,w'\sep a'=w\sep a'$.
\end{proof}

\begin{lemma}\label{lem:joinnil}
\leanref{join\_eq\_nil}
If every entry of $s$ is in $\N$ and $\join s=\eps$, then $s=\sq{}$.
\end{lemma}
\begin{proof}
If $s=\sq{w,\dots}$ is non-empty, $\join s$ begins with $w$ (it is $w$, or
$w\sep\cdots$), and $w\neq\eps$, so $\join s\neq\eps$.
\end{proof}

\begin{lemma}[Join is injective on names]\label{lem:joininj}
\leanref{join\_inj}
If every entry of $s$ and of $t$ is in $\N$ and $\join s=\join t$, then
$s=t$.
\end{lemma}

\begin{proof}
Induction on $s$, by cases on the shapes of $s$ and $t$.
\begin{itemize}
\item If $s=\sq{}$ then $\join t=\eps$, so $t=\sq{}$ by
  Lemma~\ref{lem:joinnil}. Symmetrically if $t=\sq{}$.
\item If $s=\sq{w}$ and $t=\sq{v}$ then $w=v$.
\item If $s=\sq{w}$ and $t=\sq{v,v',\dots}$, then $w=v\sep\join\sq{v',\dots}$
  contains $\sep$, contradicting $w\in\N$. Symmetrically if $s$ is the longer one.
\item If $s=\sq{w,w',\dots}$ and $t=\sq{v,v',\dots}$, then
  $v\sep\join\sq{v',\dots}=w\sep\join\sq{w',\dots}$. Since $\sep\notin w$,
  Lemma~\ref{lem:splitfirst} gives either $v=w$ and
  $\join\sq{v',\dots}=\join\sq{w',\dots}$, in which case the induction
  hypothesis gives $\sq{w',\dots}=\sq{v',\dots}$ and so $s=t$; or
  $v=w\sep a'$, which contains $\sep$, contradicting $v\in\N$. \qedhere
\end{itemize}
\end{proof}

\begin{definition}[Good shape]\label{def:shape}
\leanref{GoodShape}
A word $k$ has \emph{good shape} if it does not begin with $\sep$, does not
end with $\sep$, and does not contain $\sep\sep$.
\end{definition}

\begin{lemma}\label{lem:joinshape}
\leanref{join\_shape}
If $s\neq\sq{}$ and every entry of $s$ is in $\N$, then $\join s\neq\eps$ and
$\join s$ has good shape.
\end{lemma}

\begin{proof}
Induction on $s$. If $s=\sq{w}$, then $\join s=w$ is non-empty and contains no
$\sep$ at all, so it has good shape trivially.

If $s=\sq{w,v,\dots}$, let $J=\join\sq{v,\dots}$. By induction $J\neq\eps$
and $J$ has good shape. The word $w\sep J$ is non-empty because $w$ is.
We use Lemma~\ref{lem:splitfirst} with $\sep\notin w$ three times.
\begin{itemize}
\item \emph{No leading $\sep$.} If $w\sep J=\sep b$, that is
  $\eps\sep b=w\sep J$, then either $\eps=w$, which is false, or
  $\eps=w\sep a'$, which is impossible.
\item \emph{No trailing $\sep$.} If $w\sep J=a\sep$, that is
  $a\sep\eps=w\sep J$, then either $J=\eps$, which is false, or
  $J=a'\sep$ for some $a'$, which contradicts the good shape of $J$.
\item \emph{No $\sep\sep$.} If $w\sep J=a\sep\sep b$, that is
  $a\sep(\sep b)=w\sep J$, then either $J=\sep b$, so $J$ begins with $\sep$, or
  $J=a'\sep\sep b$. Both contradict the good shape of $J$. \qedhere
\end{itemize}
\end{proof}

% =======================================================================
\section{Chains, depth and the Key}
\label{sec:chains}

From here to the end of Section~\ref{sec:lemmas}, fix a model
$(\U,\F,\I,r,C,\nm)$.

\begin{lemma}[Chain uniqueness]\label{lem:chainunique}
\leanref{chain\_unique}
If $x_0,\dots,x_n$ and $y_0,\dots,y_m$ are chains to the same entity $x$, then
$m=n$ and $y_i=x_i$ for all $i\le n$.
\end{lemma}

\begin{proof}
Induction on $n$, for all $x$, $m$ and $y$ at once.

\emph{$n=0$.} Then $x=x_0=r$. If $m\ge1$, then $y_{m-1}\C y_m=x=r$,
contradicting AX-5. So $m=0$ and $y_0=r=x_0$.

\emph{$n=k+1$.} Then $x_k\C x$. If $m=0$, then $x=y_0=r$ and $x_k\C r$,
contradicting AX-5. So $m=j+1$ and $y_j\C x$. By AX-4, $y_j=x_k$. Now
$x_0,\dots,x_k$ and $y_0,\dots,y_j$ are both chains to $x_k$, so the induction
hypothesis gives $j=k$ and $y_i=x_i$ for $i\le k$. Hence $m=n$, and also
$y_n=x=x_n$.
\end{proof}

\begin{definition}[Depth and the chain of $x$]
\leanref{depth, chain}
By AX-6 every $x$ has a chain. By Lemma~\ref{lem:chainunique} its length and
entries do not depend on which chain is chosen. Call the length the
\emph{depth} of $x$, $\depth(x)$, and the chain \emph{the chain of $x$},
written $x_0,\dots,x_{\depth(x)}$.
\end{definition}

The Lean file picks a chain with the axiom of choice and then shows that any
other chain agrees with it (\lean{chain\_eq}). The following facts are
immediate:
\begin{itemize}
\item A prefix $x_0,\dots,x_i$ of a chain is a chain to $x_i$
  (\lean{chain\_prefix}).
\item $\depth(r)=0$, since $r$ alone is a chain to $r$, and conversely
  $\depth(x)=0$ implies $x=x_0=r$ (\lean{depth\_root},
  \lean{eq\_root\_of\_depth}).
\item No entry $x_i$ with $i\ge1$ of a chain is $r$, since $x_{i-1}\C x_i$
  and AX-5 (\lean{chain\_ne\_root}).
\end{itemize}

\begin{lemma}[One step]\label{lem:step}
\leanref{chain\_step}
If $p\C x$, then $\depth(x)=\depth(p)+1$, the chain of $x$ agrees with the
chain of $p$ at every index $i\le\depth(p)$, and its last entry is $x$.
\end{lemma}
\begin{proof}
Append $x$ to the chain of $p$. Since $p\C x$, the result is a chain to $x$
of length $\depth(p)+1$. By Lemma~\ref{lem:chainunique} it is the chain of $x$.
\end{proof}

\begin{theorem}[D1: every non-root has exactly one container]\label{thm:D1}
\leanref{D1}
If $x\neq r$, there is a $p$ with $p\C x$, and every $q$ with $q\C x$ equals $p$.
\end{theorem}
\begin{proof}
Since $x\neq r$, $\depth(x)=k+1$ for some $k$. In the chain of $x$,
$x_k\C x_{k+1}=x$, so $p=x_k$ contains $x$. Uniqueness is AX-4.
\end{proof}

\begin{definition}[Names along a chain; Key]\label{def:key}
\leanref{seqOf, Key}
For $x$ with chain $x_0,\dots,x_n$, let
\[
  \seqof(x)=\sq{\nm(x_1),\dots,\nm(x_n)},
  \qquad
  \Key(x)=\join\seqof(x) = \nm(x_1)\sep\cdots\sep\nm(x_n).
\]
The root's name is not used, and $\Key(r)=\join\sq{}=\eps$. This is
TREE-009.
\end{definition}

By construction $|\seqof(x)|=\depth(x)$ and $\seqof(r)=\sq{}$
(\lean{length\_seqOf}, \lean{seqOf\_root}). The entries of $\seqof(x)$ are names:
each $x_i$ with $i\ge1$ is not $r$, so $\nm(x_i)\in\N$ by AX-7
(\lean{seqOf\_InN}).

\begin{lemma}\label{lem:seqstep}
\leanref{seqOf\_step}
If $p\C x$, then $\seqof(x)=\seqof(p)\cat\sq{\nm(x)}$.
\end{lemma}
\begin{proof}
By Lemma~\ref{lem:step}, the chain of $x$ is the chain of $p$ followed by $x$.
\end{proof}

\begin{lemma}\label{lem:seqprefix}
\leanref{seqOf\_chain, seqOf\_take}
For the chain $x_0,\dots,x_n$ of $x$ and every $i\le n$,
$\seqof(x_i)=\pre{\seqof(x)}{i}$.
\end{lemma}
\begin{proof}
$x_0,\dots,x_i$ is a chain to $x_i$, hence by Lemma~\ref{lem:chainunique} the
chain of $x_i$. Its names are the first $i$ entries of $\seqof(x)$.
\end{proof}

\begin{lemma}[$\seqof$ is injective]\label{lem:seqinj}
\leanref{seqOf\_inj}
If $\seqof(x)=\seqof(y)$, then $x=y$.
\end{lemma}

\begin{proof}
Induction on $n=\depth(x)$. Equal sequences have equal lengths, so
$\depth(y)=n$ as well.

If $n=0$, both $x$ and $y$ are $r$.

If $n=k+1$, neither is $r$. By D1 they have containers $p\C x$ and $q\C y$,
and by Lemma~\ref{lem:seqstep}
\[
  \seqof(p)\cat\sq{\nm(x)}=\seqof(q)\cat\sq{\nm(y)}.
\]
Cancelling the last entry gives $\seqof(p)=\seqof(q)$ and $\nm(x)=\nm(y)$. Since
$\depth(p)=k$ by Lemma~\ref{lem:step}, the induction hypothesis gives $p=q$.
Then $p$ contains both $x$ and $y$, and they have one name, so $x=y$ by AX-8.
\end{proof}

This lemma is where AX-8 does its work. Every injectivity statement below
reduces to it.

% =======================================================================
\section{The lemmas L1 to L5}
\label{sec:lemmas}

\begin{theorem}[L1, chain uniqueness]
\leanref{L1}
Every $x$ has exactly one chain: a chain exists, and any two agree in length
and in every entry.
\end{theorem}
\begin{proof}
Existence is AX-6. Uniqueness is Lemma~\ref{lem:chainunique}.
\end{proof}

\begin{theorem}[L2, acyclicity]
\leanref{L2}
There are no $k\ge1$ and $y_0,\dots,y_k$ with $y_i\C y_{i+1}$ for all $i<k$
and $y_k=y_0$.
\end{theorem}

\begin{proof}
By Lemma~\ref{lem:step}, each step of $C$ raises the depth by exactly one. By
induction on $i$, $\depth(y_i)=\depth(y_0)+i$ for every $i\le k$. If $y_k=y_0$,
then $\depth(y_0)=\depth(y_0)+k$, so $k=0$, a contradiction.
\end{proof}

\begin{remark}
The research document derives L2 differently. It walks backwards from $y_0$
around the cycle by AX-4, and that walk must be the chain of $y_0$ and so
reach $r$, which then has a container. The depth argument above is
the one the Lean file uses. Both use AX-4, through
Lemma~\ref{lem:chainunique}. AX-5 and AX-6 alone do not suffice: the
structure $C=\{(r,g_1),(g_1,g_2),(g_2,g_1)\}$ satisfies both and has a cycle.
\end{remark}

\begin{theorem}[L3, Key injectivity; TREE-010]\label{thm:L3}
\leanref{L3}
If $\Key(x)=\Key(y)$, then $x=y$.
\end{theorem}
\begin{proof}
$\join\seqof(x)=\join\seqof(y)$, and all entries are names, so
$\seqof(x)=\seqof(y)$ by Lemma~\ref{lem:joininj}, and $x=y$ by
Lemma~\ref{lem:seqinj}.
\end{proof}

\begin{theorem}[L4, Key shape; TREE-011]
\leanref{L4}
$\Key(x)=\eps$ if and only if $x=r$. For $x\neq r$, $\Key(x)$ has good shape:
no leading $\sep$, no trailing $\sep$, and no $\sep\sep$.
\end{theorem}
\begin{proof}
If $x=r$, $\Key(x)=\join\sq{}=\eps$. If $\Key(x)=\eps$, then $\seqof(x)=\sq{}$
by Lemma~\ref{lem:joinnil}, so $\depth(x)=0$ and $x=r$. If $x\neq r$, then
$\seqof(x)\neq\sq{}$ (as just shown) and its entries are names, so
Lemma~\ref{lem:joinshape} applies.
\end{proof}

\begin{definition}[Proper ancestor]
\leanref{ProperAncestor}
$a$ is a \emph{proper ancestor} of $x$ if there are $k\ge1$ and
$a=z_0,z_1,\dots,z_k=x$ with $z_j\C z_{j+1}$ for every $j<k$.
\end{definition}

\begin{theorem}[L5, prefix closure; TREE-012]\label{thm:L5}
\leanref{L5}
\begin{enumerate}[(a)]
\item Let $s_1,\dots,s_n\in\N$ with $\Key(x)=s_1\sep\cdots\sep s_n$. For each
  $0<i<n$ there is a folder $a$ that is a proper ancestor of $x$ with
  $\Key(a)=s_1\sep\cdots\sep s_i$.
\item In any chain $x_0,\dots,x_n$, no $x_i$ with $i<n$ is an item.
\end{enumerate}
\end{theorem}

\begin{proof}
(a) Let $s=\sq{s_1,\dots,s_n}$. By Lemma~\ref{lem:joininj},
$\seqof(x)=s$, so $\depth(x)=n$. Let $x_0,\dots,x_n$ be the chain of $x$
and take $a=x_i$. It is a folder by AX-3, since $x_i\C x_{i+1}$. It is a proper
ancestor of $x$, witnessed by $x_i,x_{i+1},\dots,x_n$ with $n-i\ge1$ steps. And
by Lemma~\ref{lem:seqprefix},
$\Key(x_i)=\join\pre{\seqof(x)}{i}=\join\sq{s_1,\dots,s_i}$.

(b) If $i<n$ then $x_i\C x_{i+1}$, so $x_i\in\F$ by AX-3, hence $x_i\notin\I$
by AX-1 (disjoint).
\end{proof}

\begin{remark}
The hypothesis $s_j\in\N$ in (a) matters. It is what makes ``the split of
$\Key(x)$ at its separators'' well defined. Without it, $\Key(x)$ could be
written as a join of other words, for example ones containing $\sep$, whose
prefixes are not Keys.
\end{remark}

% =======================================================================
\section{L6: the trie correspondence}
\label{sec:L6}

L6 says that a model is, up to renaming its entities, nothing more than the
set of its name sequences together with a marking of which are items. For
example, the model with Keys $\eps$, \texttt{d}, \texttt{d/e} and \texttt{d/f},
where \texttt{d/f} is an item, is captured by the four sequences $\sq{}$,
$\sq{\texttt{d}}$, $\sq{\texttt{d},\texttt{e}}$ and
$\sq{\texttt{d},\texttt{f}}$, with the last one marked. To state L6 we
need to say what ``the same model'' means, and which sets of sequences are
allowed.

\begin{definition}[Isomorphism]
\leanref{Iso}
An \emph{isomorphism} between structures
\[
  S=(\U,\F,\I,r,C,\nm) \quad\text{and}\quad S'=(\U',\F',\I',r',C',\nm')
\]
is a bijection $f\colon\U\to\U'$ with
$x\in\F\iff f(x)\in\F'$, $x\in\I\iff f(x)\in\I'$, $f(r)=r'$,
$p\C x\iff f(p)\C' f(x)$, and $\nm'(f(x))=\nm(x)$.
\end{definition}

Isomorphisms compose and have inverses, and the identity is one.

\begin{definition}[Tries]
\leanref{Trie, Trie.WF}
A \emph{trie} is a pair $(L,K)$ of sets of finite sequences of words. It is
\emph{well formed} when $L$ is a tree of name sequences and $K$ marks some of
its leaves. Precisely: every entry of every sequence in $L$ is a name;
$L$ contains $\sq{}$ and is \emph{prefix closed} ($s\cat t\in L$ implies
$s\in L$); and $K$ consists of non-empty sequences in $L$, none of which has
a proper extension in $L$ (if $s\in K$ and $s\cat t\in L$, then $t=\sq{}$).
\end{definition}

Prefix closure is under dropping whole trailing \emph{names}, not trailing
characters. That is why this is a statement about sequences and not about
Key strings: in the model with Keys $\eps$, \texttt{d}, \texttt{d/f}, the
string \texttt{d/} is a character prefix of a Key and is not a Key.

\subsection{From a trie to a model}

\begin{definition}\label{def:ofTrie}
\leanref{ofTrie}
For a well-formed trie $T=(L,K)$, let $S_T$ be the structure with universe
$L$, items $K$, folders $L\setminus K$, root $\sq{}$, containment
$p\C x\iff x=p\cat\sq{w}$ for some $w$, and $\nm(s)$ the last entry of $s$
(and $\nm(\sq{})=\eps$).
\end{definition}

\begin{proposition}\label{prop:ofTrieAx}
\leanref{ofTrie\_axioms}
$S_T$ is a model of AX-1 to AX-8.
\end{proposition}

\begin{proof}
The entities are split into $K$ and $L\setminus K$, so each is exactly one
of item or folder. The root $\sq{}$ is not in $K$, since $K$ holds only
non-empty sequences, so it is a folder.

Containment appends one name. So the only possible container of $x$ is $x$
with its last entry removed, and the empty sequence has no container at
all. Every $x=\sq{w_1,\dots,w_n}$ in $L$ is reached from the root through
its prefixes $\sq{},\sq{w_1},\dots,x$. Each prefix lies in $L$ because $L$
is prefix closed, and each one contains the next.

The root's name is $\eps$, and every other sequence is named by its last
entry, which is a name. Two children $p\cat\sq{w}$ and $p\cat\sq{w'}$ of
the same $p$ with the same name have $w=w'$, so they are equal.

Finally, only folders contain anything. If $p$ contains $p\cat\sq{w}\in L$,
then $p$ has a proper extension in $L$, so $p$ is not in $K$.
\end{proof}

\begin{lemma}\label{lem:seqofTrie}
\leanref{seqOf\_ofTrie}
In $S_T$, $\seqof(s)=s$ for every $s\in L$.
\end{lemma}
\begin{proof}
The prefixes $\pre{s}{0},\dots,\pre{s}{|s|}$ form a chain to $s$, so by L1
they are the chain of $s$. The last entry of $\pre{s}{i}$ is the $i$-th
entry of $s$. So by induction on $i$, the names along the first $i$ steps are
$\pre{s}{i}$. Take $i=|s|$.
\end{proof}

\subsection{From a model to a trie}

\begin{definition}
\leanref{toTrie}
For a model $S$, let $T_S=(L_S,K_S)$ with
$L_S=\{\seqof(x):x\in\U\}$ and $K_S=\{\seqof(x):x\in\I\}$.
\end{definition}

\begin{proposition}\label{prop:toTrieWF}
\leanref{toTrie\_wf}
$T_S$ is well formed.
\end{proposition}

\begin{proof}
Every sequence in $L_S$ is the name sequence of some entity, so its entries
are names, and the root contributes $\sq{}$. Prefix closure is
Lemma~\ref{lem:seqprefix}: the prefix of length $i$ of $\seqof(x)$ is the
name sequence of $x_i$, the entity at depth $i$ on the chain of $x$. Item
sequences lie in $L_S$ by definition. They are non-empty: the root is a
folder, and no entity is both a folder and an item, so no item is the root.
Only the root has the empty sequence (L4).

The one real step is that an item's sequence has no proper extension in
$L_S$. Let $x$ be an item, and suppose $\seqof(y)$ properly extends
$\seqof(x)$. On the chain of $y$, the entity at depth $|\seqof(x)|$ has the
same name sequence as $x$, so it \emph{is} $x$ (Lemma~\ref{lem:seqinj}). It is
not the last entity of that chain, so it contains the next one, and only
folders contain anything. But $x$ is an item, and no entity is both.
\end{proof}

\begin{proposition}\label{prop:isoToTrie}
\leanref{isoToTrie}
$x\mapsto\seqof(x)$ is an isomorphism from $S$ to $S_{T_S}$.
\end{proposition}

\begin{proof}
The map is onto $L_S$ by definition, and one-to-one by
Lemma~\ref{lem:seqinj}. Its inverse picks, for each $s\in L_S$, the entity
whose sequence is $s$. It remains to see that it respects each part of the
structure.

It sends $r$ to $\sq{}$. An entity is an item exactly when its sequence is
in $K_S$, because an entity with the same sequence as an item is that item.
Folders are the remaining entities on both sides, so they correspond too.
The name of $\seqof(x)$ in $S_{T_S}$ is its last entry. That entry is
$\nm(x)$ by Lemma~\ref{lem:seqstep}, and for the root both names are $\eps$.

If $p$ contains $x$, then $\seqof(x)=\seqof(p)\cat\sq{\nm(x)}$
(Lemma~\ref{lem:seqstep}), so the images are related by containment as well.
Conversely, suppose $\seqof(x)=\seqof(p)\cat\sq{w}$. Then $x$ is not the
root, so it has a container $q$ (D1), and
$\seqof(q)\cat\sq{\nm(x)}=\seqof(p)\cat\sq{w}$. Dropping the last entry gives
$\seqof(q)=\seqof(p)$, so $q=p$ and $p$ contains $x$.
\end{proof}

\begin{lemma}\label{lem:seqiso}
\leanref{seqOf\_iso}
If $f\colon S\to S'$ is an isomorphism of models, then
$\seqof'(f(x))=\seqof(x)$ for every $x$.
\end{lemma}
\begin{proof}
Induction on $\depth(x)$. For $x=r$, $f(r)=r'$ and both sequences are
$\sq{}$. For $x\neq r$ with container $p$ of smaller depth, $f(p)\C' f(x)$,
so
$\seqof'(f(x))=\seqof'(f(p))\cat\sq{\nm'(f(x))}=\seqof(p)\cat\sq{\nm(x)}=\seqof(x)$.
\end{proof}

\subsection{The theorem}

\begin{theorem}[L6, trie correspondence]\label{thm:L6}
\leanref{L6}
\begin{enumerate}[(a)]
\item For every model $S$, $T_S$ is a well-formed trie.
\item For every well-formed trie $T$, $S_T$ is a model.
\item For every well-formed trie $T$, $T_{S_T}=T$.
\item For every model $S$, $S$ is isomorphic to $S_{T_S}$.
\item Two models $S,S'$ are isomorphic if and only if $T_S=T_{S'}$.
\end{enumerate}
\end{theorem}

\begin{proof}
(a) is Proposition~\ref{prop:toTrieWF}, (b) Proposition~\ref{prop:ofTrieAx}, and
(d) Proposition~\ref{prop:isoToTrie}.

(c) By Lemma~\ref{lem:seqofTrie}, $L_{S_T}=\{\seqof(s):s\in L\}=L$ and
$K_{S_T}=\{\seqof(s):s\in K\}=K$.

(e) If $f\colon S\to S'$ is an isomorphism, then by Lemma~\ref{lem:seqiso}
$f$ carries each $\seqof(x)$ to the same sequence, and items to items, so
$L_S=L_{S'}$ and $K_S=K_{S'}$ (the reverse inclusions use $f^{-1}$).
Conversely, if $T_S=T_{S'}$, compose the isomorphisms
$S\to S_{T_S}=S_{T_{S'}}\to S'$ from (d) (\lean{isoOfToTrieEq}).
\end{proof}

So the map $S\mapsto T_S$ is a bijection between isomorphism classes of models
and well-formed tries, with inverse $T\mapsto S_T$. The isomorphism preserves
the typing (F, I) and the names, not only the containment graph. The prose
statement of L6 leaves that open, and the Lean statement fixes it.

\begin{corollary}[Empty folders]
\leanref{L6\_empty\_folder}
In $S_T$, a leaf of $L$ that is not in $K$ is a folder that contains nothing.
\end{corollary}
\begin{proof}
It is not in $K$, so it is a folder. If it contained $x=s\cat\sq{w}\in L$,
being a leaf would force $\sq{w}=\sq{}$.
\end{proof}

% =======================================================================
\section{Infinite universes}
\label{sec:infinite}

Spec~051 adds no finiteness axiom on purpose. This section shows that this
choice has content in both directions: the axioms admit infinite models, and
with infinitely many entities, acyclicity no longer implies reachability.

\begin{theorem}[An infinite model]\label{thm:infinite}
\leanref{infinite\_model}
If $\Sigma$ has a symbol $a\neq\sep$, there is a model of AX-1 to AX-8 with
an infinite universe.
\end{theorem}

\begin{proof}
Let $L$ be the set of all finite sequences whose every entry is the one-letter
word $a$, and $K=\emptyset$. This trie is well formed. Its entries are the
name $a$, it contains $\sq{}$, and it is prefix closed. The conditions on $K$
hold because $K$ is empty. By
Proposition~\ref{prop:ofTrieAx}, $S_{(L,K)}$ is a model. The map
$n\mapsto\sq{a,\dots,a}$ ($n$ copies) is injective into its universe, since
distinct $n$ give distinct lengths.
\end{proof}

The model is a single infinite downward branch with no items, in which every
entity has finite depth. Like the one-point model $W_0$, it also witnesses
that the axioms are \emph{consistent}, so no theorem above holds merely
because nothing satisfies the hypotheses.

\begin{theorem}[Acyclic but not reachable]
\label{thm:upchain}
\leanref{infinite\_upward\_chain}
Let $\U=\{r\}\cup\Nat$, with containment $(n+1)\C n$ for every $n\in\Nat$ and
nothing else, every entity a folder, no items, and every name $\eps$. Then
every non-root has a container, $r$ has none, $C$ has no cycle, and AX-6 fails.
\end{theorem}

\begin{proof}
$n$ is contained by $n+1$, and nothing contains $r$. Along any $C$-path
$y_0,\dots,y_k$ with $k\ge1$, every $y_i$ is a number and $y_0=y_i+i$ (by
induction on $i$, as each step goes from $m+1$ to $m$). So $y_k=y_0$ forces
$k=0$. Finally, $0$ has no chain. A chain of length $0$ would make $0=r$, and
a chain of length $\ge1$ would start with $r\C x_1$, but $r$ contains nothing.
\end{proof}

All names here are $\eps$, so this structure also violates AX-7. It shows
why AX-6 cannot be replaced by acyclicity; it is not an independence witness
for AX-6 (that is row AX-6 of Table~\ref{tab:rows}).

\begin{proposition}[The finite half]\label{prop:finitehalf}
If $\U$ is finite, every non-root has a container, $r$ has none, and $C$ has
no cycle, then AX-6 holds.
\end{proposition}

\begin{proof}
Given $x$, choose a container of $x$ (if $x\neq r$), then a container of
that, and so on, as long as the current entity is not $r$. Two equal entities
in this walk would give a cycle, so its entries are distinct, and as $\U$ is
finite the walk stops. It can stop only at an entity without a container,
that is at $r$. Reversed, the walk is a chain to $x$.
\end{proof}

This proposition is \emph{not} in the Lean file. It is checked in Alloy for
universes of up to six entities (\lean{finite\_acyclic\_gives\_reach}). It
explains why, in a finite universe, acyclicity (with every non-root
contained and $r$ uncontained) already implies AX-6, so a bounded model finder
cannot separate the two. The reason the axioms keep AX-6 is
Theorem~\ref{thm:upchain}: without finiteness, acyclicity is not enough. The
other direction needs AX-4: the structure in the remark after L2 is finite,
satisfies AX-6, and has a cycle.

% =======================================================================
\section{What the bounded Alloy model adds}
\label{sec:alloy}

The Lean file proves consequences \emph{of} the axioms for every universe.
It does not show that each axiom clause is needed. That is a question about
structures that \emph{violate} the axioms, and a model finder answers it
cheaply. The Alloy model restates the eleven clauses as eleven predicates
and runs 30 commands, each with an expected outcome (a \texttt{run} must find
an instance; a \texttt{check} must find no counterexample).
\lean{scripts/alloy\_check.sh} fails if any outcome differs.

\subsection{Independence}

For each clause, Table~\ref{tab:rows} gives a finite structure that satisfies
every other clause and violates that one. Unless a row says otherwise,
$\nm(r)=\eps$ and the other entities have distinct names in $\N$. Each row is
a finite object, so checking it by hand is direct inspection. Alloy confirms
each one with a \texttt{run} at scope~4 (entities) and sequence length~6. For
the rows whose damage is a claim about Keys (AX-7 non-empty, AX-7 no-$\sep$,
AX-8), a \texttt{check} confirms that two entities get the same Key. For AX-4,
a \texttt{run} exhibits two different chains to $x$.

\begin{table}[ht]
\centering
\small
\renewcommand{\arraystretch}{1.25}
\begin{tabularx}{\textwidth}{@{}l>{\raggedright\arraybackslash}X>{\raggedright\arraybackslash}X@{}}
\toprule
Clause broken & Structure & What goes wrong \\
\midrule
AX-1 disjoint & $\U=\{r,x\}$, $\F=\{r,x\}$, $\I=\{x\}$, $C=\{(r,x)\}$ & $x$ is both folder and item. \\
AX-1 covering & $\U=\{r,x\}$, $\F=\{r\}$, $\I=\emptyset$, $C=\{(r,x)\}$ & $x$ has no kind. \\
AX-2 & $\U=\{r\}$, $\F=\emptyset$, $\I=\{r\}$, $C=\emptyset$ & The root is an item. \\
AX-3 & $\U=\{r,x,y\}$, $\F=\{r\}$, $\I=\{x,y\}$, $C=\{(r,x),(x,y)\}$ & An item contains something. \\
AX-4 & $\U=\{r,a,b,x\}$, $\F=\{r,a,b\}$, $\I=\{x\}$, $C=\{(r,a),(r,b),(a,x),(b,x)\}$ & Two chains, so two Keys, for $x$. \\
AX-5 & $\U=\F=\{r,p\}$, $C=\{(r,p),(p,r)\}$ & A cycle through $r$. \\
AX-6 & $\U=\F=\{r,g_1,g_2\}$, $C=\{(g_1,g_2),(g_2,g_1)\}$ & A detached cycle. \\
AX-7 root & $\U=\F=\{r\}$, $C=\emptyset$, $\nm(r)=\texttt{z}$ & Only a normalization breaks. No Key changes. \\
AX-7 non-empty & $\U=\{r,x\}$, $\I=\{x\}$, $C=\{(r,x)\}$, $\nm(x)=\eps$ & $\Key(x)=\eps=\Key(r)$. \\
AX-7 no $\sep$ & $\U=\F=\{r,w,a,b\}$, $C=\{(r,w),(r,a),(a,b)\}$, $\nm(w)=\texttt{a/b}$, $\nm(a)=\texttt{a}$, $\nm(b)=\texttt{b}$ & $\Key(w)=\Key(b)=\texttt{a/b}$. \\
AX-8 & $\U=\{r,x,y\}$, $\I=\{x,y\}$, $C=\{(r,x),(r,y)\}$, $\nm(x)=\nm(y)$ & $\Key(x)=\Key(y)$. \\
\bottomrule
\end{tabularx}
\caption{One countermodel per axiom clause (spec~051's rationale, \S5).
Where $\I$ is omitted, $\I=\emptyset$; where $\F$ is omitted,
$\F=\U\setminus\I$.}
\label{tab:rows}
\end{table}

The AX-7 \emph{root} row deserves a word. Nothing in the definition of $\Key$
reads $\nm(r)$, so no lemma fails without this clause. It is a normalization,
and it matters only for L6, where $\nm(r)$ would otherwise be free data that
the trie does not record.

\subsection{Consistency and witnesses}
The witness models $W_0$ (the root alone), $W_1$ (Keys $\eps$, \texttt{d},
\texttt{d/f}) and $W_2$ (Keys $\eps$, \texttt{d}, \texttt{d/e},
\texttt{d/f}, with \texttt{e} an empty folder) satisfy all eight axioms, and
their Keys are the ones listed. Together with Theorem~\ref{thm:infinite}, they
show that the axioms are consistent, in finite and in infinite universes.

\subsection{The lemmas, in bounded scope}
D1 is checked at scope~6. L1 to L5 are checked at scope~5. The Key lemmas L3 to
L5 are checked with names of at most two symbols, while L1 and L2 have no
separate bound on name length. The sequence scope is chosen so that no structure in scope is
excluded by the encoding of Keys. With at most five entities the depth is at
most~4, and a Key has at most one name and one $\sep$ per level and no
trailing $\sep$, so its length is at most $4\cdot3-1=11$. A separate
\texttt{run} shows that length~11 is actually reached.

Alloy cannot state L6. L6 quantifies over pairs of structures and over
isomorphisms between them, and an Alloy instance is a single structure. It
also cannot reach Theorems~\ref{thm:infinite} and~\ref{thm:upchain}. The checks are
therefore evidence for finite universes up to the stated scope, and the Lean
proofs are what extend them to all universes.

\subsection{The checks are not vacuous}
A check that passes because the axioms are too strong proves nothing. The
model was mutated by deleting one conjunct at a time from the axioms
(mutants run once and not committed). Without AX-4, D1, L1 and L2 fail.
Without AX-8, L3 fails. Without AX-3, L5 fails. Without AX-7 (no~$\sep$),
L3, L4 and L5 fail. This matches the proofs above. AX-4 enters directly in D1 and through
Lemma~\ref{lem:chainunique}, AX-8 through Lemma~\ref{lem:seqinj}, AX-3 through
Theorem~\ref{thm:L5}, and the no-$\sep$ clause through
Lemma~\ref{lem:joininj} (L3, L5) and Lemma~\ref{lem:joinshape} (L4).

\subsection{Division of labour}
\begin{center}
\small
\begin{tabular}{@{}lll@{}}
\toprule
Question & Alloy (bounded) & Lean (all universes) \\
\midrule
Each clause independent & yes, 11 rows & --- \\
Axioms consistent & $W_0$, $W_1$, $W_2$ & infinite model, Thm.~\ref{thm:infinite} \\
D1, L1 to L5 & up to scope 5--6 & proved \\
L6 & cannot be stated & proved \\
Finite half of \S\ref{sec:infinite} & up to scope 6 & --- \\
Infinite upward chain & out of reach & proved \\
\bottomrule
\end{tabular}
\end{center}

% =======================================================================
\appendix
\section{Concordance with the sources}

\begin{center}
\small
\begin{tabular}{@{}lll@{}}
\toprule
Here & Spec 051 & Lean declaration \\
\midrule
AX-1 to AX-8 & TREE-001 to TREE-008 & \lean{Axioms} (fields \lean{ax1\_disjoint} \dots\ \lean{ax8}) \\
Lemma~\ref{lem:splitfirst} & --- & \lean{split\_first} \\
Lemma~\ref{lem:joininj} & --- & \lean{join\_inj} \\
Lemma~\ref{lem:joinshape} & --- & \lean{join\_shape} \\
Lemma~\ref{lem:chainunique} & --- & \lean{chain\_unique} \\
Lemma~\ref{lem:step} & --- & \lean{chain\_step} \\
Theorem~\ref{thm:D1} (D1) & --- & \lean{D1} \\
Definition~\ref{def:key} ($\Key$) & TREE-009 & \lean{Key} \\
Lemma~\ref{lem:seqinj} & --- & \lean{seqOf\_inj} \\
L1, L2 & --- & \lean{L1}, \lean{L2} \\
L3 & TREE-010 & \lean{L3} \\
L4 & TREE-011 & \lean{L4} \\
L5 & TREE-012 & \lean{L5} \\
Theorem~\ref{thm:L6} (L6) & --- & \lean{L6}, \lean{L6\_empty\_folder} \\
Theorem~\ref{thm:infinite} & --- & \lean{infinite\_model} \\
Theorem~\ref{thm:upchain} & --- & \lean{infinite\_upward\_chain} \\
Proposition~\ref{prop:finitehalf} & --- & (Alloy only) \\
\bottomrule
\end{tabular}
\end{center}

The Lean file uses only Lean's core library and its standard axioms
(propositional extensionality, choice, quotient soundness), never
\lean{sorry}. A \lean{\#guard\_msgs} block pins the axioms each headline result
depends on. Choice enters where this text says ``the chain of $x$'' and
``the entity whose sequence is $s$'', and in the excluded middle used for
AX-1 (covering) of $S_T$ (Proposition~\ref{prop:ofTrieAx}). The last is why
Theorem~\ref{thm:infinite} depends on choice although it uses no chain.

\end{document}
